\documentclass[a4paper,12pt]{article}

\usepackage[utf8]{inputenc}
\usepackage[T1]{fontenc}
\usepackage[francais]{babel}
\usepackage{mathptmx}
\usepackage[scaled]{helvet}
\usepackage{courier}

\usepackage[margin=2cm,nohead,foot=1.4cm]{geometry}
\usepackage{parskip}
\usepackage{comment}
\usepackage{graphicx}
\usepackage{fancyvrb}

\begin{document}

\begin{flushleft}
\begin{tabular}[t]{l}
\Large\textsf{\textbf{Rapport sur le manuscrit de Thèse présenté par}}\\[4mm]
\Large\textsf{\textbf{Jonathan Pastor}}\\[4mm]
\large\textsf{\textbf{pour l'obtention du Doctorat de l'École des Mines de Nantes}}\\
\end{tabular}%

\vspace{4mm}

\noindent%
\begin{tabular}[t]{l}
Rapporteur: Stéphane \textsc{Genaud},\\
Professeur des Universités,\\
ENSIIE.\\[5mm]
Titre du document évalué :\\\\[3mm]
\end{tabular}%

\begin{quote}
\Large \textit{\textsf{\textbf{
Contributions à la mise en place d'une infrastructure de Cloud Computing à large
échelle}}}
\end{quote}

\end{flushleft}

\vspace{3mm}

Le document  présenté parJonathan  Pastor décrit les  recherches qu'il  a menées
dans le cadre de sa thèse sous la direction conjointe de Adrien Lèbre, chargé de
recherches à l'INRIA Rennes  Bretagne-Atlantique, et Frédéric Desprez, directeur
de recherches à l'INRIA Grenoble Rhône-Alpes. 



\subsection*{Contexte de l'étude}

Les  infrastructures  de  Cloud  Computing,  dans  lesquelles  des  fournisseurs
hébergent des services informatiques, sont  ajourd'hui hébergées dans des centre
de données dont  la taille ne cesse  de croître afin de  favoriser des économies
d'échelles. Une  alternative envisageable  à cette  tendance de  centraliser les
ressources, est  au contraire de décentraliser  ces centres de données  vers des
infrastrucures  géographiquement  réparties,  comme  celles  dont  dispoent  les
opérateurs  de  réseaux, afin  de  rapprocher  les ressources  des  utilisateurs
finaux. Cependant,  les difficultés  qu'engendrent la gestion  décentralisée des
ressources  sont  considérables.  Le  travail  de la  thèse  s'inscrit  dans  ce
contexte, en proposant des solutions logicielles aux problèmes posés.



\subsection*{Analyse du document présenté}

Le document est organisé en deux parties comptant 8 chapitres hors les
conclusions et perspectives.

Après une introduction  globale de la structure du document,  la première partie
(II) présente le  contexte de la thèse. Après une  brève introduction historique
de  l'évolution de  l'informatique jusqu'au  cloud computing  (chapitre 1),  une
description assez  complète du  paysage du cloud  computing est  faite (chapitre
2). Le chapitre  suivant se focalise sur la description  et comparaison de trois
architectures  possibles pour  les  centres de  données et  le  modèle de  cloud
computing induits : méga-centres de données,  fédération de centre de données et
edge/fog computing. Le  dernier chapitre d'analyse du  contexte actuel (chapitre
4) traite  du logiciel  OpenStack largement utilisé  aujourd'hui pour  gérer les
clouds de  type IaaS.  Ce  chapitre devient  plus technique et  décrit utilement
l'architecture  du   logiciel  pour  engager   une  réflexion  sur   les  moyens
envisageables pour  décentraliser son  fonctionnement. Différentes  approches de
déploiement  hiérarchique,  publiés  dans  la littérature,  sont  présentées  de
manière  assez  synthétique.   Toutefois,   une  description  plus  précise  des
difficultés  liées aux  mécanismes  de synchronisation  et d'exclusion  mutuelle
auraient pu y trouver leur place (par  exemple, page 64). Au final, cette partie
descriptive  du  contexte  est  très  pédagogique et  démontre  une  très  bonne
connaissance du domaine par l'auteur.

La  deuxième partie  (III)  aborde  les contributions  propres  du travail.   Le
chapitre 5  pose clairement  les objectifs d'un  système Locality  based Utility
Computing (LUC-OS)  qui répondrait à l'objectif  général de la thèse  de pouvoir
assurer la gestion  de ressources décentralisées. Ce chapitre est  court, car il
ne présente  que les  fonctionnalités désirées et  le choix  d'adapter OpenStack
pour y répondre. Deux éléments  d'OpenStack sont identifiés comme problématiques
pour  son  adaptation à  un  environnement  distribué  :  le bus  de  messagerie
(RabbitMQ qui utilise un agent de mise en relation centralisé) et son système de
base  de données  (MySQL). L'obstacle  du bus  de messagerie  est écarté  car la
communauté s'est emparée du problème, et  c'est l'obstacle de la base de données
centralisée qui est au coeur du travail  et fait l'objet du chapitre suivant. 

Le chapitre 6 (\emph{Support des bases  clé/valeur dans Nova}) est beaucoup plus
conséquent. La proposition faite ici consiste à substituer un système de base de
données de type clé/valeur au système de base de données existant dans le module
Nova  d'OpenStack. Il  faut souligner  que tenter  d'insérer ses  idées dans  un
logiciel  répandu, en  constante  évolution et  d'une  complexité certaine,  est
ambitieux et requiert une approche méticuleuse.  Le chapitre explique clairement
les choix  possibles d'un  point de  vue technique  et leurs  implications, pour
justifier le choix fait de développer  une interface ORM nommée Rome, permettant
à  OpenStack d'accéder  au système  de stockage  de type  clé/valeur Redis.   Le
chapitre détaille  ensuite les développements  nécessaires dans l'ORM  Rome pour
que  les  opérations relationnelles  disponibles  dans  l'API, soient  également
fonctionnelles quand  le moteur de  gestion des  données est Redis.   Cet aspect
couvre  les jointures,  les sessions  qui doivent  faire appel  à un  système de
verrou  distribué, et  les optimisations  des requêtes  vis à  vis du  volume de
données  transférées.   La  fin  du  chapitre est  consacrée  à  une  évaluation
expérimentale      de     cette      réalisation     sur      la     plate-forme
Grid'5000. L'expérimentation en  environnement réel doit être  saluée, car c'est
une étape chronophage mais cruciale  pour la vérification des hypothèses faites,
et  malheureusement  pas  toujours  présente  dans  les  travaux  menés  par  la
communauté. Sur  la forme,  la présentation  du protocole  expérimental pourrait
être  améliorée  concernant  la  description  des  matériels  impliqués  et  les
caractéristiques  des  machines  virtuelles  (VMs) créées.   La  description  de
l'expérience n'est ici  pas assez complète pour en  assurer sa reproductibilité,
ce qui  est paradoxal quand  le document fait  état de l'utilisation  de l'outil
Execo  qui  a  précisément   cet  objectif.   Néanmoins,  l'expérience  apparaît
pertinente,  comparant  les comportements  du  prototype  dans un  environnement
constitué  d'un   seul  cluster,   puis  dans  des   configurations  multi-sites
(artificiellement créés  par émulation  d'une augmentation  de la  latence). Ces
résultats  expérimentaux  montrent que  le  prototype  utilisant un  système  de
stockage clé/valeur permet  la création de VMs dans un  temps constant jusqu'à 8
sites, alors que  la version originale avec base de  données relationnelles voit
un brusque quadruplement de temps de création en passant de 6 à 8 sites.

Au final,  ce chapitre  décrit une  réelle contribution,  dont le  prototype est
implanté dans le  logiciel phare du domaine. La démarche  expérimentale menée en
environnement réel renforce grandement la  proposition, même si une analyse plus
fine des résultats me semble souhaitable.


Le chapitre 7 (\emph{Prise en compte de la localité réseau: le cas de DVMS}) est
la deuxième  contribution de la  thèse. L'algorithme DVMS, précédent  travail de
l'équipe,  est  un algorithme  distribué  destiné  à calculer  dynamiquement  le
placement des machines  virtuelles afin d'atteindre des  critères d'équilibre de
charge.   Cet  algorithme utilise  des  mécanismes  pair-à-pair pour  structurer
l'ensemble des  machines virtuelles  en groupes  logiques (overlay  structuré de
Chord).  Cette structuration ne prend pas  en compte la localité réseau (définie
par la latence)  des machines et le  travail de thèse a  pour objectif d'étendre
DVMS  pour minimiser  les  collaborations  entre noeuds  lointains  en terme  de
latence réseau.  Le  chapitre commence par citer les efforts  analogues faits au
début des  années 2000 dans la  recherche sur les réseaux  pair-à-pair.  Ensuite
est développé l'algorithme  proprement dit, dont l'objectif  est le regroupement
de machines virtuelles proches du point de vue réseau. L'information de localité
réseau est obtenu ici grâce à Vivaldi.

Quelques  points  d'ombres   me  semblent  subsister  dans   la  description  de
l'algorithme  concernant la  terminaison  et les  enrôlements concurrents.   Une
description  formelle de  l'algorithme eût  été la  bienvenue.  Il  faut relever
l'effort  d'intégration logicielle  fait ensuite  en intégrant  le mécanisme  de
structuration,  Chord  ou Vivaldi,  comme  un  module  spécialisé de  DVMS.   Le
chapitre s'achève  avec une validation  expérimentale sur Grid'5000,  avec cette
fois,  une  configuration  entièrement   réelle,  distribuée  sur  quatre  sites
géographiquement distants. DVMS calcule dynamiquement  le placement des VMs pour
satisfaire des critères  d'équilibrage de charge, déclenchant  des migrations de
VMs pour  maintenir la  satisfaction des  contraintes.  L'expérience  observe le
nombre de  migrations intra-site et  inter-sites qui  ont lieu, avec  la version
DVMS basée sur Chord  et celle sur Vivaldi. La version  Vivaldi tenant compte de
la localité  réseau remplit  parfaitement l'objectif avec  un taux  de migration
intra-site environ deux fois supérieure à la version originale. L'expérience est
convaincante,  bien que  le  seul  jeu de  paramètres  choisi pour  l'expérience
concernant  la distribution  statistique de  la charge  n'est pas  discuté alors
qu'il correspond peut être à des situtations courantes.


Le chapitre 8 (\emph{Contributions à VMPlaceS : un simulateur pour algorithme de
  placement  de machines  virtuelles}) est  la troisième  contribution de  cette
thèse.   Le logiciel  VMPlaceS est  un simulateur  s'appuyant sur  le moteur  de
simulation SimGrid, bien  connu de la communauté. L'objectif est  de simuler les
résultats d'algorithmes  de placements de  machines virtuelles, étant  donné une
charge  de travail  évolutive dans  le  temps et  une infrastructure  modélisée.
L'orientiation  du travail  vers la  simulation  dans cette  dernière partie  du
travail est  pertinente.  Les  deux chapitres précédents  faisaient appel  à une
expérimentation  en  environnement  réel,  à échelle  réduite,  et  le  document
jsutifie bien  que la  simulation est un  moyen complémentaire  d'expérimenter à
plus grande échelle. La démarche expérimentale  suivie dans ce chapitre est tout
à fait  recommandable :  les résultats  de la simulation  de VMPlaceS  sont tout
d'abord comparés à une exécution réelle sur une infrastructure réelle (réduite à
32 noeuds) afin de valider le comportement du simulateur. Ensuite, la validation
à plus grande échelle est faite en simulation.

La comparaison expérimentale porte ici  sur le comportement de trois algorithmes
de calcul de placements des VMs  : l'un centralisé, le deuxième hiérarchique, et
le troisième  distribué (DVMS  présenté au chapitre  précédent). La  métrique de
comparaison est le temps cumulé où les VMs n'ont pas pu satisfaire la contrainte
de  charge  maximale  imposée,  c'est-à-dire qu'elles  sont  en  surcharge.  Les
résultats montrent une nette réduction du temps de surcharge avec DVMS, validant
ainsi l'approche.



\subsection*{Conclusion}

Le  document est  bien rédigé  et agréable  à lire.   Jonathan Pastor  décrit le
domaine de  recherche dans  lequel il a  inscrit ses travaux  avec un  recul qui
démontre ses connaissances poussées des  techniques et technologies de l'état de
l'art. Une  partie importante du  travail est à la  mise en {\oe}uvre  des idées
présentées dans  des logiciels existants,  en particulier le  logiciel OpenStack
très répandu.   C'est un  choix ambitieux,  qui a pour  avantage de  vérifier la
faisabilité de  idées dans  un contexte  proche de  celui des  environnements de
production, mais  le désavantage  d'hériter de toute  la complexité  du logiciel
utilisé,  rendant plus  difficile la  compréhension fine  du prototype  proposé.
Pour valider les objectifs visés,  une évaluation expérimentale considérable est
menée  pour chacune  des  propositions. Les  résultats expérimentaux  confirment
l'intérêt des contributions proposées.

Les  travaux ont  fait l'objet  de  plusieurs publications  (dans 3  conférences
internationales de très  bon niveau et un chapitre de  livre).  Considérant tous
ces éléments, j'émets un avis favorable à  la soutenance de la thèse de Jonathan
Pastor en vue de l'obtention du titre  de Docteur en informatique de l'École des
Mines de Nantes. La thèse peut être soutenue en l'état.

\vspace{.8cm}
\begin{flushright}
Stéphane \textsc{Genaud},\\
{\small fait à Strasbourg, le 20 septembre 2016}.\\
\includegraphics[width=.20\textwidth]{signgenaud.jpg}
\end{flushright}
\end{document}
